\documentclass[12pt,a4paper]{article}
\usepackage[margin=1in]{geometry}
\usepackage{natbib}
\usepackage{bibentry}
\bibliographystyle{apalike}
\pagestyle{empty}

\begin{document}
\nobibliography*
\section*{Jingxian Liu}

\noindent\textbf{Fields/Subfields:} Labor Economics, Demographic Economics, Applied Econometrics

\medskip
\noindent\textbf{Keywords:} Wage inequality, Intergenerational mobility,
Demographic structure, Population aging, Fertility, Gender inequality,
Labor supply, Household consumption

\renewcommand{\refname}{Papers}
\nocite{autor2008,card2013,kleven2019,chetty2014,jaimovich2009}
\bibliography{bib/references}

\section*{The R.A.P. Framework}

\bigskip
\noindent\bibentry{autor2008}.

\begin{description}
  \item[Research Question:] What explains the sharp rise in U.S. wage inequality since the 1980s?
  \item[Answer:] The rise in wage inequality was not simply a one-time event, and these patterns may be partly explained by a broader version of the skill-biased technological change hypothesis.
  \item[Positioning:] Some recent studies argue that wage inequality rose mainly in the early 1980s because of the declining real minimum wage and changes in the workforce. This paper revisits that argument by looking separately at the upper and lower parts of the wage distribution and considering the role of longer-term changes in the demand for skills.
\end{description}

\bigskip
\noindent\bibentry{card2013}.

\begin{description}
  \item[Research Question:] How do changes in employer wage premiums and worker–employer matching help explain the rise in wage inequality in West Germany?
  \item[Answer:] Rising dispersion in worker-specific and establishment-specific wage components, together with stronger sorting of high-wage workers into high-premium establishments, accounts for nearly all of the increase in wage inequality.
  \item[Positioning:] Earlier explanations focused on changes in skill supply and demand, while studies of employer wage premiums offered little evidence on how their contribution to inequality changed over time. This paper distinguishes changes in worker heterogeneity, establishment premiums, and worker sorting to examine how much workplaces contributed to rising wage inequality.
\end{description}

\bigskip
\noindent\bibentry{kleven2019}.

\begin{description}
  \item[Research Question:] Why does gender earnings inequality persist despite the narrowing of gender gaps in education and the introduction of anti-discrimination policies?
  \item[Answer:] Female child penalties account for most of the remaining gender earnings inequality.
  \item[Positioning:] Persistent gender earnings gaps despite convergence in education motivate examining parenthood beyond traditional decompositions that condition on occupation and other job characteristics. Building on earlier studies of motherhood penalties, this paper treats those characteristics as potential channels and quantifies how child penalties contribute to aggregate gender inequality over time.
\end{description}

\bigskip
\noindent\bibentry{chetty2014}.

\begin{description}
  \item[Research Question:] What does intergenerational mobility look like in the United States as a whole and across its regions?
  \item[Answer:] Intergenerational mobility varies across areas, with higher upward mobility in areas with less residential segregation, less income inequality, better primary schools, greater social capital, and greater family stability.
  \item[Positioning:] Earlier research mainly compared relative mobility across countries, but differences in data and methods made comparisons difficult and shed little light on local opportunities for children from low-income families. This paper uses a common administrative dataset to measure relative and absolute mobility across the United States and examine the local characteristics associated with each.
\end{description}

\bigskip
\noindent\bibentry{jaimovich2009}.

\begin{description}
  \item[Research Question:] How do changes in the age composition of the workforce affect business cycle volatility in the G7 countries?
  \item[Answer:] Changes in workforce age composition account for a significant fraction of the variation in business cycle volatility across the G7 and approximately one-fifth to one-third of the moderation in U.S. volatility.
  \item[Positioning:] Research on the Great Moderation focused on the decline in U.S. volatility after the mid-1980s, while earlier studies showed that young and older workers were more sensitive to business cycles. This paper links these age differences to demographic shifts across the G7 to help explain both the earlier rise and later decline in U.S. volatility.
\end{description}

\end{document}
